\scanpage{84}
\mspage{77}
% archive sheet number: 83
7. Filtrations cardinales et parties accessibles d'une \unsure{gross.|grosse} catégorie.

Je me réfère à la théorie des foncteurs et catégories accessibles de SGA 4, I, \S\ 9 (p. 138--178). Je me donne une grosse catégorie (disons une $\unsure{\mathcal U|U}$-catégorie) $E$ avec une \uline{filtration cardinale} (croissante)
\[
(\mathrm{Filt}^{\pi}(E))_{\pi\geq\pi_0},
\]
\margin{$\pi_0$ un card. infini fixé}
par des sous-catégories str. pleines $\mathrm{Filt}^{\pi}(E)$.

On \unsure{a|.}

a) $\mathrm{Filt}^{\pi}(E)$ \deleted{\unsure{est}} essentiellement petites, i.e. \unsure{équivalentes|équivalente} à une petite catégorie (\unsure{ou.|ou} à une catégorie $\in\mathcal U$)

b) $E$ satisfait $L_{\pi_0}$ (i.e. $E$ est stable par $\varinjlim$ de type $I$, quand $I$ est un ens. ordonné grand devant $\pi_0$), et $\mathrm{Filt}^{\pi}(E)$ est stable par les limites inductives filtrantes indexées par $I$ grand devant $\pi_0$ et tel que $\operatorname{card}I\leq\pi$.

c) Pour tt $\pi\geq\pi_0$ et pour tt $X\in\mathrm{Ob}\,E$, on a
\[
X\simeq\varinjlim_{i\in I}X_i
\]
où les $X_i$ ds $\mathrm{Filt}^{\pi}(E)$, et où $I$ est un ens. ordonné grand devant $\pi$.
% slanted
\margin{NB b) et c) impliquent que $E=\bigcup_{\pi\geq\pi_0}\mathrm{Filt}^{\pi}(E)$, \deleted{\uline{NB} \unsure{On a donc}}}
De plus, si $X\in\mathrm{Ob}\,\mathrm{Filt}^{\pi'}(E)$, avec $\pi'\unsure{\geq|>}\pi$, on peut prendre $I$ tel que $\operatorname{card}I\leq\pi'^{\pi}$ (donc si $\pi'=2^{c'}$ \inserted{avec $c'\geq\pi$}, d'où $(\pi')^{\unsure{n|\pi}}=(2^{c'})^{\pi}=2^{c'\pi}=2^{c'}=\pi'$, on aura $\operatorname{card}I\leq\pi'$. (\uline{NB} Si on admet l'hyp. du continu généralisée, on aura tjrs \deleted{$\operatorname{card}I\leq$} $\pi'^{\pi}=\pi'$ donc $\operatorname{card}I\leq\pi'$)
